\section{Генерация параметров с помощью утилиты \texttt{ecgen}}

Для подготовки входных данных в работе использовалась внешняя утилита \texttt{ecgen}~\cite{ecgen}. Она предназначена для генерации параметров эллиптических кривых и позволяет задавать свойства генерируемой кривой через набор флагов командной строки. В рамках настоящей работы применялись только те возможности утилиты, которые были нужны для построения тестовых примеров над простыми полями.

Для генерации кривых над простым полем $\mathbb{F}_p$ использовался ключ \texttt{--fp}, после которого указывается битовая длина модуля поля. Например, команда вида
\begin{verbatim}
ecgen --fp 192
\end{verbatim}
задаёт поиск кривой над простым полем размерности 192 бита.

Для получения кривой с простым порядком подгруппы использовался флаг \texttt{-p} (длинная форма --- \texttt{--prime}). Такой режим соответствует криптографически корректному базовому сценарию, когда стойкость ECDLP должна определяться большой подгруппой простого порядка.

Для построения примеров, демонстрирующих уязвимость групп с ``гладким'' порядком, применялся флаг \texttt{--smooth=BOUND}. Он задаёт верхнюю границу на битовую длину простых множителей порядка группы.

Для генерации аномальных кривых, уязвимых к атаке Смарта, использовался флаг \texttt{--anomalous}.

Команды, использовавшиеся при подготовке тестовых данных, имели вид
\begin{minted}{bash}
# кривые без особенностей
./ecgen --fp -c1 -p -r 40 > dataset/general/40bit.json
./ecgen --fp -c3 -p -r 32 > dataset/general/32bit.json
./ecgen --fp -c5 -p -r 16 > dataset/general/16bit.json

# кривые с составным порядом для атаки Pohlig-Hellman
./ecgen --fp -c5 -r --smooth=32 128 > dataset/smooth/128bit32smooth.json
./ecgen --fp -c5 -r --smooth=32 64 > dataset/smooth/64bit32smooth.json
./ecgen --fp -c5 -r --smooth=16 64 > dataset/smooth/64bit16smooth.json

# кривые с простым порядком, совпадающим с порядком поля для Smart Attack
./ecgen --fp -c5 -r --anomalous 512 > dataset/anomalous/512bit.json
./ecgen --fp -c5 -r --anomalous 256 > dataset/anomalous/256bit.json
./ecgen --fp -c5 -r --anomalous 128 > dataset/anomalous/128bit.json
\end{minted}

Вывод \texttt{ecgen} имеет формат JSON. В простейшем случае он содержит описание поля, коэффициенты уравнения кривой, порядок группы точек и список найденных подгрупп. Типичная структура результата имеет вид
\begin{minted}{json}
{
    "field": {
        "p": "0x..."
    },
    "a": "0x...",
    "b": "0x...",
    "order": "0x...",
    "subgroups": [
        {
            "x": "0x...",
            "y": "0x...",
            "order": "0x...",
            "cofactor": "0x...",
            "points": [
                "x": "0x...",
                "y": "0x...",
                "order": "0x..."
            ]
        }
    ]
}
\end{minted}

Поле \texttt{field.p} задаёт модуль простого поля, поля \texttt{a} и \texttt{b} определяют коэффициенты формы Вейерштрасса, а поле \texttt{order} задаёт число точек на кривой. Массив \texttt{subgroups} содержит информацию о выделенной подгруппе: координаты порождающей точки, её порядок и кофактор (отличие порядка подгруппы от порядка группы).